% ============================================================================
%  Section III — Dataset and Measurement Setup  (IEEE TIM, IEEEtran)
%  Passive voice / third person throughout.
%  Requires: booktabs, graphicx, multirow, tikz.
%  Details, and the reused measured-speed values, are taken from the source
%  data article \cite{ertargin2026smartphone}; the dataset is on Mendeley Data
%  \cite{ertargin2026dataset} (CC BY 4.0).
% ============================================================================

\section{Dataset and Measurement Setup}
\label{sec:dataset}

The present study is based on the publicly available smartphone-acquired
vibration dataset introduced in~\cite{ertargin2026smartphone} and archived on
Mendeley Data~\cite{ertargin2026dataset}. The dataset was selected because it
embodies exactly the low-cost, low-bandwidth acquisition regime targeted in this
work. This section summarises the test rig, the acquisition chain, the emulated
fault conditions, and the organisation of the recordings; the reader is referred
to the original data article for the complete experimental protocol. The
acquisition and dataset parameters are consolidated in
Table~\ref{tab:specs}.

\subsection{Test Rig and Machine}
The measurements were carried out on a $1.1$~kW three-phase squirrel-cage
induction machine, a rating representative of the large population of small
general-purpose drives found in industry. Variable-speed operation was obtained
through a voltage/frequency ($V/f$)-controlled inverter, and a mechanical load
was applied by a DC generator coupled directly to the motor shaft, enabling both
loaded and unloaded operation. A consumer smartphone was affixed to the motor
terminal box with double-sided adhesive tape, oriented approximately parallel to
the shaft axis; the mounting location was marked so that the placement could be
reproduced across successive installations. The complete arrangement of motor,
inverter, and DC generator is shown in Fig.~\ref{fig:testrig}.

\subsection{Acquisition Chain and Signal Channels}
Vibration signals were acquired with an iPhone~15 Pro Max using its integrated
tri-axial MEMS inertial sensor, accessed through the \textit{Sensor Play}
application. Six channels were recorded for every operating condition: the three
raw acceleration components ($g_x$, $g_y$, $g_z$) and the three
gravity-compensated linear-acceleration components ($g_{userx}$, $g_{usery}$,
$g_{userz}$) produced by the device's on-board sensor fusion. All channels were
sampled at $100$~Hz, which corresponds to the maximum stable rate supported by
the acquisition application, and each operating condition was recorded
continuously for $15$~minutes. The resulting sampling rate is one to three
orders of magnitude below that of the industrial-grade benchmarks discussed in
Section~\ref{sec:related}, and constrains the analysable band to below
$50$~Hz in accordance with the Nyquist criterion~\cite{nyquist1924certain,
shannon1948mathematical}.

\subsection{Health States and Fault Emulation}
Six motor health states were considered: one healthy (fault-free) baseline and
five fault conditions spanning mechanical and electrical origins, as summarised
in Table~\ref{tab:faults}. The bearing faults were introduced by substituting a
$6205$ZZ bearing that had been extracted from industrially operated motors, so
that the degradation reflects realistic in-service mechanisms rather than
artificial machining. Two lubrication-starvation severities were distinguished:
insufficient lubrication~(B1) and severe insufficient lubrication~(B2), the
severity being assigned qualitatively from the rotational behaviour observed
during manual inspection; a bearing with a cracked outer ring~(B3) constituted a
localised mechanical defect. The electrical fault, voltage imbalance~(V), was
produced by inserting voltage-divider resistors in one supply phase, which
yielded a pronounced phase-current asymmetry (for example, phase currents of
$0.75$/$2.75$/$3.03$~A at $50$~Hz under load~\cite{ertargin2026smartphone}). The
rotor fault, broken rotor bar~(R), was emulated by drilling four holes of
$5$--$9$~mm diameter and $16$--$18$~mm depth into the rotor bars. Because
fault-related vibration signatures are speed-dependent, the shaft speed was
recorded for every condition; the measured values are reproduced in
Table~\ref{tab:speeds}, where the reduced speed under voltage imbalance
(e.g., $1326$~r/min versus $1434$~r/min for the healthy motor at $50$~Hz under
load) is consistent with the torque disturbance introduced by that fault.

\subsection{Operating Conditions and Dataset Organisation}
Each of the six health states was recorded at three supply frequencies
($30$, $40$, and $50$~Hz) and under two load conditions (loaded and unloaded),
producing a fully crossed matrix of $6\times3\times2 = 36$ recordings. The files
follow the naming convention \texttt{Class\_Load\_SpeedHz}, where the load
indicator is~\texttt{1} for loaded and~\texttt{0} for unloaded operation; for
instance, \texttt{B1\_1\_40Hz} denotes insufficient lubrication under load at
$40$~Hz. Each recording is provided in both comma-separated (\texttt{.csv}) and
MATLAB (\texttt{.mat}) formats. Representative gravity-compensated waveforms for
the six health states, recorded under identical operating conditions, are shown
in Fig.~\ref{fig:signals}; the healthy motor exhibits markedly lower vibration
amplitude, whereas the fault conditions produce larger and morphologically
distinct responses.

\subsection{Data Integrity}
Prior to analysis, an independent integrity check was performed on all $36$
recordings. Each file was confirmed to contain exactly $89{,}999$ samples per
channel, corresponding to $900$~s at $100$~Hz, with no missing values and no
systematic sample-and-hold duplication (the maximum fraction of consecutive
duplicated rows in any recording was below $0.01\%$). The channel set and
recording length were therefore uniform across the entire dataset, and no
imputation or resampling was required.

% ---------------------------------------------------------------------------
\begin{table}[t]
\centering
\caption{Acquisition and dataset specifications.}
\label{tab:specs}
\renewcommand{\arraystretch}{1.2}
\begin{tabular}{@{}ll@{}}
\toprule
\textbf{Parameter} & \textbf{Value} \\
\midrule
Motor & $1.1$~kW three-phase squirrel-cage IM \\
Speed control & $V/f$-controlled inverter \\
Loading & DC generator (shaft-coupled) \\
Sensor & iPhone~15 Pro Max, tri-axial MEMS \\
Acquisition app & Sensor Play \\
Channels & $6$ ($g_x,g_y,g_z$; $g_{userx},g_{usery},g_{userz}$) \\
Sampling rate & $100$~Hz \\
Record length & $15$~min ($89{,}999$ samples/channel) \\
Health states & $6$ (H, B1, B2, B3, V, R) \\
Supply frequencies & $30$, $40$, $50$~Hz \\
Load conditions & loaded, unloaded \\
Recordings & $36$ ($6\times3\times2$) \\
Formats & \texttt{.csv}, \texttt{.mat} \\
Repository & Mendeley Data, CC BY 4.0 \\
\bottomrule
\end{tabular}
\end{table}

% ---------------------------------------------------------------------------
\begin{table}[t]
\centering
\caption{Health states and fault-emulation summary.}
\label{tab:faults}
\renewcommand{\arraystretch}{1.25}
\begin{tabular}{@{}clp{4.2cm}@{}}
\toprule
\textbf{Code} & \textbf{Health state} & \textbf{Emulation} \\
\midrule
H  & Healthy & Fault-free baseline \\
B1 & Insufficient lubrication & Reduced lubricant, $6205$ZZ bearing; severity set by manual inspection \\
B2 & Severe insufficient lubrication & Advanced lubricant starvation; severity set by manual inspection \\
B3 & Cracked outer ring & Bearing with cracked outer ring (ex-service degradation) \\
V  & Voltage imbalance & Voltage-divider resistors in one supply phase; phase-current asymmetry \\
R  & Broken rotor bar & Four holes drilled in rotor bars ($5$--$9$~mm diameter, $16$--$18$~mm depth) \\
\bottomrule
\end{tabular}
\end{table}

% ---------------------------------------------------------------------------
\begin{table}[t]
\centering
\caption{Measured shaft speed (r/min) for each health state, supply frequency,
and load condition (reproduced from~\cite{ertargin2026smartphone}).}
\label{tab:speeds}
\renewcommand{\arraystretch}{1.15}
\setlength{\tabcolsep}{4pt}
\begin{tabular}{@{}llcccccc@{}}
\toprule
\textbf{Supply} & \textbf{Load} & \textbf{H} & \textbf{B1} & \textbf{B2} &
\textbf{B3} & \textbf{V} & \textbf{R} \\
\midrule
\multirow{2}{*}{$30$~Hz} & Loaded   & 833.9 & 829.4 & 833.0 & 833.8 & 787.0 & 815.8 \\
                         & Unloaded & 892.7 & 888.3 & 892.7 & 896.3 & 893.0 & 894.9 \\
\multirow{2}{*}{$40$~Hz} & Loaded   & 1136 & 1138 & 1138 & 1140 & 1076 & 1122 \\
                         & Unloaded & 1190 & 1192 & 1193 & 1192 & 1192 & 1192 \\
\multirow{2}{*}{$50$~Hz} & Loaded   & 1434 & 1436 & 1439 & 1436 & 1326 & 1425 \\
                         & Unloaded & 1486 & 1490 & 1491 & 1489 & 1486 & 1486 \\
\bottomrule
\end{tabular}
\end{table}

% ---------------------------------------------------------------------------
% Test-rig photograph with numbered TikZ callouts overlaid on the included image.
% Save the photograph as fig_testrig.jpg next to this file. The callout
% coordinates are normalised to the image (x: 0=left..1=right, y: 0=bottom..1=top);
% nudge the (x,y) pairs below if a marker does not sit on its component.
\begin{figure}[t]
\centering
\begin{tikzpicture}
  \node[anchor=south west,inner sep=0] (img) at (0,0)
      {\includegraphics[width=\columnwidth]{fig_testrig.jpg}};
  \begin{scope}[x={(img.south east)},y={(img.north west)}]
    \tikzset{callout/.style={circle,draw=black,line width=0.4pt,fill=white,
             font=\bfseries\footnotesize,inner sep=0.5pt,minimum size=1.35em}}
    \node[callout] at (0.90,0.32) {1};   % induction machine
    \node[callout] at (0.64,0.33) {2};   % DC generator (load)
    \node[callout] at (0.77,0.28) {3};   % shaft coupling
    \node[callout] at (0.80,0.54) {4};   % smartphone (MEMS sensor)
    \node[callout] at (0.66,0.62) {5};   % V/f variable-frequency supply
    \node[callout] at (0.85,0.90) {6};   % oscilloscope
    \node[callout] at (0.30,0.45) {7};   % resistive load bank
    \node[callout] at (0.52,0.30) {8};   % digital multimeter
  \end{scope}
\end{tikzpicture}
\caption{Test-rig arrangement (photograph adapted from~\cite{ertargin2026smartphone}):
(1)~squirrel-cage induction machine under test; (2)~shaft-coupled DC generator
providing the mechanical load; (3)~flexible shaft coupling; (4)~smartphone
supplying the tri-axial MEMS vibration signal; (5)~$V/f$ variable-frequency
supply; (6)~oscilloscope; (7)~resistive load bank; (8)~digital multimeter.}
\label{fig:testrig}
\end{figure}

% ---------------------------------------------------------------------------
\begin{figure}[t]
\centering
\includegraphics[width=\columnwidth]{fig_example_signals.png}
\caption{Representative gravity-compensated vibration ($g_{userx}$, $3$~s excerpt)
for the six health states under identical operating conditions (loaded,
$50$~Hz supply). Amplitudes share a common scale.}
\label{fig:signals}
\end{figure}
